\section{Computação}

A documentação pública identifica componentes de software, mas não detalha toda a eletrônica. A apresentação conjunta de 2023 descreve a evolução de instalações manuais e acesso remoto por scripts para serviços em contêineres e Kubernetes~\cite{oliveira2023tevel}. A Tabela~\ref{tab:tecnologias_documentadas} reúne registros de diferentes momentos e contextos.

\begin{table*}[!htb]
    \centering
    \caption{Tecnologias documentadas, funções e limites da evidência.}
    \label{tab:tecnologias_documentadas}
    \small
    \setlength{\tabcolsep}{5pt}
    \renewcommand{\arraystretch}{1.10}

    \begin{tabularx}{\textwidth}{@{}
        >{\raggedright\arraybackslash}p{0.27\textwidth}
        >{\raggedright\arraybackslash}p{0.35\textwidth}
        >{\raggedright\arraybackslash}X
        @{}}
        \hline
        \textbf{Tecnologia e fonte} &
        \textbf{Função} &
        \textbf{Limite} \\
        \hline

        GPU na estação terrestre~\cite{reeuwijk2025} &
        Processamento de visão computacional. &
        Modelo, quantidade e consumo não informados. \\
        \hline

        K3s, Flux e NVIDIA Operator~\cite{oliveira2023tevel} &
        Orquestração, sincronização de aplicações e suporte à GPU. &
        Controlador de voo não identificado. \\
        \hline

        Palette e Kairos~\cite{oliveira2023tevel} &
        Gestão e atualizações com particionamento A/B. &
        Confiabilidade da missão não demonstrada. \\
        \hline

        TPM e criptografia~\cite{oliveira2023tevel} &
        Proteção de dados persistentes. &
        Segurança integral não demonstrada. \\
        \hline

    \end{tabularx}
\end{table*}

A presença de Kubernetes não comprova a execução das malhas de estabilização nesse ambiente. As fontes analisadas não especificam o modelo do controlador de voo, as taxas das malhas, o escalonamento ou o isolamento entre controle e missão. Na análise de \acp{VANT}, interromper um serviço de gestão não equivale, por si só, a perder a autoridade sobre os atuadores.